\documentclass{article}
\usepackage[utf8]{inputenc}
\usepackage{geometry}
\usepackage{xcolor}
\geometry{a4paper, margin=1in}

\title{Evaluación de Respuestas: Gabriel}
\author{Evaluador Prof. CESAR RODRIGUEZ}
\date{\today}

\begin{document}

\maketitle

\section*{Análisis del Cuestionario}

A continuación se presenta la evaluación de tus respuestas a la "Evaluación de la Semana 1: 
Consultas DNS (A/AAAA)".

\subsection*{Pregunta 1: Diferencia entre registros A y AAAA}
\begin{itemize}
    \item \textbf{Tu Respuesta (transcrita):} \textit{"Registros A: IPV4. Registros AAAA: IPV6."}
    \item \textbf{Evaluación:} \textbf{\textcolor{green}{Acertada}}.
    \item \textbf{Comentarios:} Tu respuesta es concisa y precisa. Captura correctamente la diferencia fundamental entre ambos tipos de registro, asociando cada uno a su respectiva versión del Protocolo de Internet.
\end{itemize}

\subsection*{Pregunta 2: Propósito de \texttt{dns.resolver.resolve()}}
\begin{itemize}
    \item \textbf{Tu Respuesta (transcrita):} \textit{"Recibe una url a consultar y según los parametros en el Array muestra los registros soncitados."}
    \item \textbf{Evaluación:} \textbf{\textcolor{orange}{Regular}}.
    \item \textbf{Comentarios:} Aunque la idea general es correcta, la terminología es imprecisa. La función no recibe una "URL", sino un nombre de dominio. Su propósito principal es "resolver" un dominio para un tipo de registro específico y devolver los resultados, no solo "mostrarlos". La respuesta demuestra una comprensión funcional pero carece de la precisión técnica esperada.
\end{itemize}

\subsection*{Pregunta 3: Múltiples registros 'A'}
\begin{itemize}
    \item \textbf{Tu Respuesta (transcrita):} \textit{"ESTO es para balanceo de carga y alta disponibilidad. El dns ira rotando entre las direcciones disponibles."}
    \item \textbf{Evaluación:} \textbf{\textcolor{green}{Acertada}}.
    \item \textbf{Comentarios:} Excelente. Identificas correctamente que esta práctica se utiliza para el balanceo de carga y la alta disponibilidad. Además, describes acertadamente el mecanismo de "rotación" (conocido como DNS Round Robin), que distribuye el tráfico entre múltiples servidores.
\end{itemize}

\section*{Conclusión de la Evaluación}

Tu desempeño en la evaluación es muy bueno. Demuestras una sólida comprensión de los conceptos de DNS relacionados con los registros A/AAAA y las técnicas de balanceo de carga. La única área de mejora es la precisión en la terminología técnica al describir funciones de bibliotecas de software.

\begin{itemize}
    \item \textbf{Pregunta 1:} Acertada
    \item \textbf{Pregunta 2:} Regular
    \item \textbf{Pregunta 3:} Acertada
\end{itemize}

\textbf{CALIFICACIÓN: 8.33 en base a 10 puntos.}

\end{document}